% Generated draft: requires mathematical and source review.
\documentclass[11pt]{article}
\usepackage[margin=25mm]{geometry}
\usepackage{fontspec}
\setmainfont{Latin Modern Roman}
\usepackage{amsmath,amssymb,mathtools,unicode-math}
\setmathfont{Latin Modern Math}
\usepackage{enumitem,needspace,xurl,hyperref}
\setlength{\parindent}{0pt}
\setlength{\parskip}{5pt}
\setlength{\emergencystretch}{3em}
\newcommand{\N}{\mathbb{N}}
\newcommand{\Z}{\mathbb{Z}}
\newcommand{\Q}{\mathbb{Q}}
\newcommand{\R}{\mathbb{R}}
\newcommand{\C}{\mathbb{C}}
\newcommand{\F}{\mathbb{F}}
\newcommand{\A}{\mathbb{A}}
\newcommand{\PP}{\mathbb{P}}
\newcommand{\E}{\mathbb{E}}
\newcommand{\eps}{\varepsilon}
\newcommand{\dd}{\,\mathrm{d}}
\begin{document}
\section*{UnsolvedMath ID: 30006630}
\textbf{Source problem code:} catalogue entry
\subsection{Definitions and mathematical statement}
For \(s\geq 0\), let \(H^s(\R)\) denote the Sobolev space
\[
H^s(\R)=\left\{f\in\mathcal S'(\R):\widehat f\text{ is a function and }\int_{\R}(1+|\xi|^2)^s|\widehat f(\xi)|^2\,\dd\xi<\infty\right\}.
\]
In particular, \(H^1(\R)\) consists of square-integrable functions with square-integrable weak derivative. Set
\[
I=\left(\frac13,\frac12\right),\qquad J=\left(\frac37,\frac12\right).
\]
The supplied record concerns the mass-critical generalized Korteweg--de Vries equation on \(\R\). Its original excerpt refers to equation (1) and Theorem 1, but neither is included in full. Consequently, the equation's normalization and sign, the precise definition of the power-rate exponent \(\nu\), and the hypotheses of the constructed family cannot be recovered from the supplied text. The questions below retain those source-dependent terms explicitly; they are not a complete specification of a solution class.

The intended blowup residue is the remainder obtained after subtracting the singular bubble from a solution and approaching its blowup time. The excerpt uses the time orientation \(t\downarrow0\). It does not supply the bubble, its modulation parameters, or the topology in which the remainder converges. Thus, saying that the residue belongs to \(H^1(\R)\) must be distinguished from asserting convergence in \(H^1(\R)\), which is not specified here.

\begin{enumerate}
\item For every solution in the source's intended class that blows up in finite time at a finite spatial point, with power-rate exponent \(\nu\in I\), does a blowup residue exist and belong to \(H^1(\R)\)? The supplied assessment says that this is open even for the particular solution attributed to reference [6] at \(\nu=2/5\). That reference is not identified in the supplied excerpt.
\item For each \(\nu\in J\), do the corresponding finite-time, finite-point blowup solutions attributed to Theorem 1 enjoy any form of stability? The source does not specify a perturbation space, a topology, permitted changes in blowup parameters, or a precise stability property.
\item Determine whether the only \(H^1\) power blowup rates are \(\nu=2/5\) and \(\nu>3/7\). Restricted to the finite-time, finite-point range \(I\), the proposed classification can be expressed as
\[
\mathcal R_{H^1}\cap I\stackrel{?}{=}\left\{\frac25\right\}\cup\left(\frac37,\frac12\right),
\]
where \(\mathcal R_{H^1}\) denotes the set of power-rate exponents realized by solutions in the source's intended \(H^1\) class. This notation is conditional on specifying that solution class and its rate convention. The supplied phrase ``\(H^1\) power blowup rates'' does not explicitly distinguish regularity of the solution from regularity of its residue, and the scope of the classification outside \(I\) is not specified. Within \(I\), the displayed candidate also excludes the endpoint \(\nu=3/7\).
\item If the proposed classification holds, construct, for every
\[
\nu\in\left(\frac13,\frac37\right)\setminus\left\{\frac25\right\},
\]
a finite-time, finite-point blowup solution with that power-rate exponent and infinite energy, belonging to a suitable \(H^\kappa(\R)\) with \(0<\kappa<1\). The record does not specify the energy functional or what ``infinite energy'' means, nor whether one common \(\kappa\) must work for all such \(\nu\) or \(\kappa\) may depend on \(\nu\).
\end{enumerate}

The excerpt additionally attributes an \(H^1(\R)\) residue \(r_\alpha\) to the construction, with \(r_\alpha(x)=0\) for \(x<0\) and an asserted local power profile for \(0<x<\delta\), near the blowup point placed at the origin. Its displayed power, the coefficient involving \(c_\alpha\) and \(Q\), and the relation between \(\alpha\) and \(\nu\) are corrupted by text extraction. The excerpt also ends mid-sentence. These expressions cannot be imposed as precise formulas from the supplied evidence. The corrected statement supplies \(\nu=2/5\); the corresponding fraction in the original excerpt is corrupted. These source defects require mathematical review before a fully specified formulation can be accepted.
\subsection{Short English statement}
The questions ask whether finite-time blowup at a finite point leaves an \(H^1\) remainder, whether the constructed families are stable, and which power blowup rates are possible in \(H^1\). Conditional on the proposed classification, the remaining specified rates are to be constructed at suitable regularity \(H^\kappa\), \(0<\kappa<1\), with infinite energy. Several essential definitions and formulas are missing or corrupted in the supplied excerpt.
\subsection{Sources}
\begin{itemize}
\item Local UnsolvedMath registry, permanent ID \(30006630\), source problem code \texttt{catalogue entry}, title ``Blowup Residues, Stability, and Rates for Mass-Critical gKdV,'' in the Oberwolfach Reports Open Problems collection, proposed year 2026. No individual proposer is supplied.
\item Supplied report citation: \emph{Nonlinear Waves and Dispersive Equations} (2026), Oberwolfach Reports, DOI \texttt{10.4171/owr/2026/10}, pp.\ 611--682. \url{https://doi.org/10.4171/owr/2026/10}.
\item Supplied literature source, ``OWR report'': \url{https://publications.mfo.de/bitstream/handle/mfo/4444/OWR_2026_10.pdf?sequence=1}.
\item Supplied literature source, ``Martel--Pilod preprint'': \url{https://arxiv.org/abs/2601.20801}. No title or complete bibliographic entry is supplied, and its identification with the excerpt's reference [6] is not established by the record.
\item The supplied literature assessment was checked on \texttt{2026-08-22} and labels the problem open. It reports that both existence of the requested \(H^1\) residue and stability of the constructed solutions remain open, even at \(\nu=2/5\). This is supplied dated curation, not an independently verified literature assessment; the record itself requests independent verification of time-sensitive status and cited claims.
\item The registry labels the statement \texttt{corrected\_verified} and says that it was reconstructed and checked against the official Oberwolfach report. This is supplied provenance, not independent mathematical verification of this draft. The supplied original excerpt references equation (1), Theorem 1, Remark 5, and reference [6], with missing context and corrupted formulas.
\end{itemize}
\end{document}
